\scanpage{46}
catégoriques généraux \ldots\ ].

Donc en ce cas, l'invariant suffit à réconstituer \deleted{\unsure{toute}} la situation $(X,X',s)$.

Dans le cas où $X$ est normal, intègre, et où \uline{$i$ n'est pas identiquement égal à $0$ ou $1$}, \unsure{comme on connaît} le \unsure{revêt} non ramifié de $X$ sur l'ouvert $U$ où $i$ et $i-1$ inversibles $[S_{i(i-1)}]$, on le normalise partout, d'où reconstitution de $(X,X')$; d'autre part, on connaît $s$ sur $U$, donc ($X$ normal) partout.

Donc ici encore, \uline{$i$ permet de tt reconstituer}. Ici cependant, \unsure{se pose la} question : quels $i$ conviennent ??

Il faut exprimer (i) qu'une certaine extension de corps -- de degré $3$ -- du corps des fonctions de $X$ est non ramifié sur $X$ (ii) Qu'une certaine application rationnelle de $X$ qui s'en déduit est partout définie \ldots
\mspage{45}
